\section{Własna niesprzeczność i granice zaufania: Gödel II oraz Löb}\label{sec:granice-zaufania}
\subsection*{System nie wystawia sam sobie dostatecznego certyfikatu}\label{sec:17:1}
Drugie twierdzenie dotyczy precyzyjnego zdania $\operatorname{Con}_T$: „nie istnieje kod dowodu sprzeczności w $T$”, przy ustalonym standardowym kodowaniu. Gdy $T\supseteq\mathrm{PA}$ jest niesprzeczna i jej aksjomaty można efektywnie wyliczać, zdanie to nie ma dowodu w $T$. W tym skrypcie nie udowodniliśmy analogicznej tezy dla każdej słabszej teorii ani dla dowolnego zdania potocznie nazwanego „niesprzecznością”.

Jeżeli ufamy tylko regułom i aksjomatom $T$, to nie uzyskamy z nich własnego standardowego zapewnienia, że żadna kombinacja tych reguł i aksjomatów nigdy nie doprowadzi do sprzeczności.

Nie wynika stąd, że dowód niesprzeczności jest niemożliwy gdziekolwiek. Silniejsza teoria $U$ może dowodzić $\operatorname{Con}_T$. Wtedy uzasadnienie opiera się na założeniach $U$, a nie wyłącznie na $T$. Jeśli do $U$ stosuje się to samo drugie twierdzenie, analogiczne ograniczenie pojawia się dla jej własnego $\operatorname{Con}_U$. Takie względne dowody nadal mają wartość: pokazują dokładnie, jakie założenia wystarczają do uzasadnienia danej teorii.

\subsection*{Każdy sprawdzony przypadek to nie to samo co wszystkie naraz}\label{sec:17:2}
Jeśli $T$ jest niesprzeczna, każdy konkretny standardowy numer $n$ nie jest kodem dowodu sprzeczności. Lemat~\ref{res:7:7} daje osobno dla każdego $n$:
\[
T\vdash\neg\operatorname{Prf}_T(\overline n,\ulcorner\bot\urcorner).
\]
Można też połączyć dowody dla dowolnie długiego, ale ustalonego skończonego odcinka. Nie daje to automatycznie jednego dowodu zdania $\forall p\,\neg\operatorname{Prf}_T(p,\ulcorner\bot\urcorner)$. To ostatnie zdanie jest właśnie $\operatorname{Con}_T$. Przy założeniach drugiego twierdzenia zachodzi $T\nvdash\operatorname{Con}_T$: wszystkie konkretne sprawdzenia nie składają się na jeden dostępny w $T$ dowód ogólnej niesprzeczności.

Różnica przypomina rozdział między sprawdzeniem dowolnej podanej liczby przypadków a posiadaniem jednego argumentu obejmującego wszystkie przypadki. W arytmetyce wolno stosować indukcję, ale trzeba jeszcze mieć dowód właściwego kroku indukcyjnego; nie dostaje się go za darmo z samej rodziny poprawnych sprawdzeń.

\begin{przyklad}[Niesprzeczna teoria może błędnie twierdzić, że istnieje dowód sprzeczności]\label{res:17:1}
Załóżmy, że PA jest niesprzeczna. Wtedy teoria $U=\mathrm{PA}+\neg\operatorname{Con}_{\mathrm{PA}}$ jest również niesprzeczna. Gdyby bowiem była sprzeczna, twierdzenie o dedukcji dawałoby $\mathrm{PA}\vdash\neg\neg\operatorname{Con}_{\mathrm{PA}}$, a klasyczna logika $\mathrm{PA}\vdash\operatorname{Con}_{\mathrm{PA}}$, wbrew Twierdzeniu~\ref{res:12:2}.

Jednocześnie dodany aksjomat $\neg\operatorname{Con}_{\mathrm{PA}}$ jest przy naszym założeniu fałszywy w standardowym $\mathbb N$. Głosi istnienie kodu dowodu sprzeczności w PA, choć rzeczywistego skończonego dowodu nie ma. Każdy model $U$ musi więc mieć niestandardowego świadka tego zdania: żaden standardowy numerał nie może nim być, bo $Q$ poprawnie obala każdego takiego konkretnego kandydata.

Nie otrzymaliśmy rzeczywistego dowodu sprzeczności w PA. Otrzymaliśmy model, który uznaje pewien niestandardowy element za kod takiego dowodu. To przykład różnicy między niesprzecznością a $\Sigma_1$-poprawnością, nie obalenie drugiego twierdzenia.
\end{przyklad}

\subsection*{Co naprawdę mówi twierdzenie Löba}\label{sec:17:3}
Zdanie $\Box\varphi\to\varphi$ brzmi: „jeśli $T$ dowodzi $\varphi$, to $\varphi$ jest prawdziwe”, gdzie po prawej nie ma słowa „prawdziwe”, lecz samo zdanie $\varphi$. Twierdzenie Löba mówi, że jeśli wewnątrz $T$ da się udowodnić taką implikację dla konkretnego $\varphi$, to $T$ już dowodzi $\varphi$. Wymagane są tu warunki wyprowadzalności (D1)--(D3).

Nie można więc uzyskać nowego zdania przez łatwą sztuczkę: „nie potrafię go dowieść, ale udowodnię, że mój system na pewno nie dowodzi go błędnie”. W tej formalnej postaci udowodnienie takiej gwarancji jest już wystarczająco silne, aby otrzymać samo zdanie. Dla $\varphi=\bot$ gwarancja staje się niesprzecznością, dlatego twierdzenie Löba zawiera drugie twierdzenie jako szczególny przypadek.

Nie znaczy to, że $T$ nie dowodzi żadnej implikacji $\Box\varphi\to\varphi$. Jeśli $T\vdash\varphi$, taka implikacja wynika od razu z logiki. Ograniczenie dotyczy zdań, których $T$ jeszcze nie dowodzi, i wewnętrznego uzasadniania refleksji. Zewnętrzne zaufanie do prawdziwości aksjomatów nie jest tym samym co dowód odpowiedniej zasady w samym systemie.

\subsection*{Dlaczego (D1)--(D3) nie obiecują prawdziwości twierdzeń}\label{sec:17:4}
Warunki Hilberta--Bernaysa--Löba opisują operacje na dowodach: można potwierdzić istnienie konkretnego dowodu, połączyć dowody przez modus ponens oraz potwierdzić dowodliwość stwierdzenia o dowodliwości. Są to gwarancje składniowe. Żaden z tych warunków nie jest ogólną zasadą $\Box\varphi\to\varphi$.

Podobnie program sprawdzający formalny dowód może poprawnie ustalić, że wniosek wynika z zadanych aksjomatów, bez rozstrzygania, czy wszystkie te aksjomaty są prawdziwe w zamierzonym świecie. Sprawdzanie poprawności kroków, poszukiwanie dowodu i uzasadnianie prawdziwości aksjomatów to trzy odrębne zadania.

